\newpage
\phantomsection
\label{prf:computational-geometry-linear-functionals-convex-combinations}
\LRAProofFor{lem:computational-geometry-linear-functionals-convex-combinations}

\begin{remark*}[Return]
\hyperref[lem:computational-geometry-linear-functionals-convex-combinations]{Return to Lemma}
\end{remark*}
\ProofVaultURL{proof-vault/pending/computational-geometry-linear-functionals-convex-combinations}

\begin{theorem*}[Linear functionals preserve convex combinations]
Let \(d\in\mathbb{R}^2\), and let
\[
  z=\sum_{i=1}^{n}\lambda_i q_i
\]
be a convex combination of points \(q_1,\dots,q_n\in\mathbb{R}^2\). Then
\[
  \langle d,z\rangle
  =
  \sum_{i=1}^{n}\lambda_i\langle d,q_i\rangle.
\]
Consequently, if \(\langle d,q_i\rangle\le M\) for every \(i\), then
\[
  \langle d,z\rangle\le M.
\]
\end{theorem*}

\begin{proof}
\textbf{Professional Standard Proof.}~\\
TODO: Use linearity of the dot product in the second variable and the
nonnegative weights summing to \(1\).
\end{proof}

\begin{proof}
\textbf{Detailed Learning Proof.}~\\
TODO: Expand the computation of
\(\langle d,\sum_i \lambda_i q_i\rangle\), then derive the inequality from
\(\langle d,q_i\rangle\le M\) and \(\sum_i\lambda_i=1\).
\end{proof}

\begin{remark*}[Proof structure]
The proof is algebraic: distribute the linear functional across the finite sum,
then use the convex-combination weights.
\end{remark*}

\begin{dependencies}
\begin{itemize}
  \item \hyperref[def:computational-geometry-convex-combination]{Convex Combination}
  \item TODO: dot product / linear functional
\end{itemize}
\end{dependencies}

\clearpage
